\chapter{把二十六章压成一页}
\chaptermeta{I}

\begin{ledetext}
合上书两年之后，细节会忘。这一页让你在五分钟内，把二十六章重新拼回一条链。
\end{ledetext}

这一页是这本书的替身。

你合上书两年，细节会掉，最后只剩几个模糊的念头；到那天再翻回二十六章，成本太高，也太慢。所以这里把全书压成一张五分钟能扫完的图：一条链，四卷，二十六句，最后五件事。忘得只剩这一页，也够你把整座建筑重新立起来。

用法只有一种：忘了，就翻到这一页扫一遍；扫完还拿不准，顺着那句话翻回原文。

\section{一条链}

全书只讲这一条链，反复讲，换着角度讲。它一共十二环。

\begin{chainflow}
\chainnode{钱是记账}\chainarrow \chainnode{记账要有人信}\chainarrow \chainnode{信任靠制度和抵押品}\chainarrow \chainnode{银行放贷时创造存款}\chainarrow \chainnode{钱约等于债}\chainarrow \chainnode{债要付息}\chainarrow \chainnode{付息需要增长}\chainarrow \chainnode{增长需要投资}\chainarrow \chainnode{投资需要给风险定价}\chainarrow \chainnode{定价错误加杠杆等于危机}\chainarrow \chainnode{危机后重置资产负债表}\chainarrow \chainnode{循环重来}
\end{chainflow}

十二环连起来读，是一口气：钱是记账，记账要有人信，信任靠制度和抵押品垫着；所以银行放贷那一刻，存款是当场造出来的，钱约等于债。债要付息，付息就得有增长，增长靠投资，投资得先给风险定价。定价一旦错了，再叠上杠杆，就是危机；危机逼人重置资产负债表，重置完，下一轮从头再来。

这条链上还有一个位置，常被看漏：你没有站在外面看它，你站在末端，是一张很小的资产负债表。机器怎么转你左右不了，但那张小表怎么摆，是你唯一说了算的事。这也是四卷为什么一路从天上讲到地上。

这条链还能倒着读。一场危机过后，钱没有消失，只是换了名字记在别人账上；从“重置资产负债表”往回摸，能摸到谁在借、谁在还、谁在扛损失。

随便拿一条新闻试。看到“某银行债券利差走阔”，往链条上找位置：这是“投资需要给风险定价”那一环，市场在重估谁会还不上钱。看到“居民提前还贷”，那是“钱约等于债”这一环在收缩。每一条新闻都能挂到其中一环上。

\section{四卷各管哪一段}

四卷不是按主题分的，是按这条链的进度分的。前一卷故意留的坑，后一卷负责填。

地基篇管头三环：钱是记账，记账要有人信，信任靠制度和抵押品。它给的不是知识，是一套读账的语法：资产、负债、净值，利息、通胀、贴现。解决的是“黑话太多”这件事，让没学过经济的人先能听懂饭桌上的对话。它回答“钱凭什么值钱”，不回答“怎么变有钱”；还故意留了一个坑：银行放出去的钱，最初从哪儿冒出来。

管道篇接过那个问题，管第四到第八环：放贷创造存款，钱约等于债，债要付息，付息需要增长，增长需要投资。它把货币怎么被造出来、央行管什么、利率怎么传到你的月供、债股汇怎么标价，一节一节铺开，让新闻标题背后的管道变得能看穿。走到结尾，你已经能把一条新闻翻译成一段因果。它留下一个更大的问号：管道里流的东西，为什么每隔几年就炸一次。

天气篇管第九到第十一环：给风险定价、加杠杆、出危机、重置资产负债表。它讲清楚三件事：风险的价格怎么算，杠杆怎么把小麻烦放大成灾难，泡沫和周期为什么必然反复。解决的是“为什么好日子和坏日子轮流来”。读完你会明白，危机不是意外，是这台机器的固定零件。它不管一件事：这些潮汐落到你个人头上，该怎么办。

落地篇管最后半环，也就是你：怎么做自己的资产负债表，收益到底从哪来，怎么配置才不白担风险，复利哪句是真哪句是假，骗局靠什么零件运转，2026 年哪六个变量在改写算式。从这里往后，书不再教你去看懂别人，只教你摆平自己那张表。它不再往下留问题，因为它绕回了起点：钱是记账，记账要人信，而你此刻已经知道谁在欠你。

四卷可以跳着读。急着防骗，直接翻落地篇；想懂新闻，先补管道篇。但那条链，最好从头到尾各走一遍。

\section{二十六章，压成二十六句}

下面每句都是那一章真正的落点，不是标题的复述。按顺序读，是全书骨架；随手挑一句，那一章就该在眼前立起来。这二十六句也不是并列的，是把上面那条链拆成二十六段；哪一句看着眼生，就说明那一环你已经忘了。

\begin{flowlist}
  \flowitem{00}{一张一百块的环球旅行}{你花掉的 98 块没动过一毫米，只是两行数字一减一加，却已顺手参与了一次清算、一次汇率报价。}
  \flowitem{01}{钱不是财富，钱是一套记账系统}{先有欠账，后有铸币；钱不是藏在金库里的东西，是一群人一起记住的账。}
  \flowitem{02}{你的存款，其实不在银行的金库里}{你把钱存进银行的那一刻，它就不再属于你，剩下的只是银行欠你的一张条。}
  \flowitem{03}{利息：时间是有价格的}{利息拆开看，是时间、通胀、违约、流动性四样东西的价，反过来就是贴现率。}
  \flowitem{04}{通货膨胀：每天从你口袋里拿走一点的那个人}{通胀不改动任何一行数字，它只把量东西的那把尺子悄悄改短，再在债主和欠债人之间搬钱。}
  \flowitem{05}{资产、负债、净值：全世界通用的三个词}{资产等于负债加净值，净值才是财富；而一个人的金融资产，必然是另一个人的负债。}
  \flowitem{06}{所有金融的原点：今天的钱换明天的钱}{金融只干两件事：把今天的钱搬到明天，把风险从一个人摊给一群人，其余都是工程细节。}
  \flowitem{07}{货币是怎么“凭空”出现的}{银行放贷的那一刻钱才出生，你还款的那一刻这笔钱就死了，从没有人问过一句“这能兑现吗”。}
  \flowitem{08}{央行不是印钞厂，是水闸和温度计}{央行决定不了钱有多少，它定的是钱的价格；它能把马牵到河边，逼不了马喝水。}
  \flowitem{09}{利率的传导链：从政策利率到你的房贷}{政策利率到你的月供之间隔着七道闸，每一道都可能把信号拖慢，甚至让它蒸发在半路。}
  \flowitem{10}{债券：世界上最大、最无聊、最要命的市场}{债券的价格和收益率永远反向，期限越长跌得越狠，因为那是一个除法，不是行情。}
  \flowitem{11}{股票：把一家公司切成一亿片}{股票是排在所有人后面的那份剩饭，风险最大，也正因如此，长期回报最高。}
  \flowitem{12}{汇率与美元体系：为什么全世界用别人的钱做生意}{汇率是两种钱互相的报价，今天主要看资本往哪儿跑，货往哪儿走是次要的。}
  \flowitem{13}{风险的价格：为什么有人借钱 3\%，有人 12\%}{利率是一叠拆得开的账单，看不清的那部分按最坏的算，抵押品和征信能把它压下来。}
  \flowitem{14}{杠杆：金融世界的放大镜与绞肉机}{杠杆不改变你的对错，只改变你能撑多久；放大收益和让你归零的是同一个数。}
  \flowitem{15}{衍生品：从农民的保险单到华尔街的赌桌}{衍生品是切下来转手的风险，它是不是赌，取决于你身上原本有没有那份方向相反的风险。}
  \flowitem{16}{影子银行：管道之外的管道}{影子银行做的是同一门生意，拿短钱做长事，只是出事时没有存款保险，也没有最后贷款人。}
  \flowitem{17}{泡沫与崩溃的通用剧本}{泡沫要四个零件：真实的新故事、便宜的信贷、杠杆、新进场的人；故事通常是真的，错的是价格和时间表。}
  \flowitem{18}{周期：为什么好日子和坏日子会轮流来}{周期不是钟表，是四五个不同长度的钟摆挂在一起；你能判断站在哪一段，算不出下一个拐点。}
  \flowitem{19}{给自己做一张资产负债表}{净值是算出来的，不是感觉出来的；你表上最大的一笔资产是你自己，而它没上保险。}
  \flowitem{20}{收益从哪儿来：三种钱和一个残酷等式}{你赚的每一块钱只有三个来路：经济增长、别人的口袋、被摩擦吃掉，后两种早晚找上你。}
  \flowitem{21}{资产配置：唯一的免费午餐}{把不同步的东西放在一起，收益不减、波动变小，它的开关是相关性，不是数量。}
  \flowitem{22}{复利的真相，和关于复利的谎言}{复利是乘法，怕中间有个特别小的数，更怕中间有个 0；所以先别出局，再谈收益。}
  \flowitem{23}{骗局的五个零件}{骗局只有五个零件：看不懂的故事、又高又稳的承诺、紧迫、熟人、退出摩擦，认全了外壳随便换。}
  \flowitem{24}{2026：六个你必须知道的新变量}{六件事同时在变天：AI 烧钱、支付重做、黄金上位、私募膨胀、利率上移、被动资金变巨兽。}
  \flowitem{25}{金融的本质，是陌生人之间的信任}{金融是让陌生人敢签字的机器，技术换了六次，三个问题没换：谁承诺、凭什么信、出事谁兜底。}
\end{flowlist}

书里每个数字都标了 2026 年 8 月。两年后回来看，数字已经不对，但这一页没有一句依赖某个数字。数字当水位读，机制当河道读，河道不搬家。

\begin{pullquote}{}
大多数人研究钱的顺序是反的：先问买什么，最后才问钱从哪儿来。而骗局和危机，吃的都是这个顺序。
\end{pullquote}

\section{如果只记住五件事}

这五条不是排在最前面的五章，是全书里最值钱的五个结论。它们的共同点：不需要任何背景，能直接压在一个具体决定上。

\begin{takeaway}{如果只记住五件事}
从上面二十六句里再蒸馏出五条。忘掉其余全部，也不会被骗，也能看懂新闻。

\begin{itemize}
  \item 任何“钱”都是某个人的负债。先问谁欠你、拿什么撑着，再谈别的。
  \item 钱是被银行借出来的，不是印出来的。大家愿意借，钱才变多；抢着还钱，钱就变少。
  \item 利率是全世界的除数。它一动，股票、房子、债券、你的月供一起重算。
  \item 高收益是风险的价格，不是善意。“又高又稳”四个字凑齐，基本是骗局。
  \item 你只有一条时间线。先保证不出局，再谈收益；储蓄率是唯一全归你管的变量。
\end{itemize}

\end{takeaway}

这五条不用记原文，记意思就够。真拿不准一笔钱，回到第一条：谁欠你，拿什么撑。

回到序章那个扫码的晚上。你扫了 98 块，老陈账上多 98，你账上少 98，什么都没变，又什么都变了：这一回，你看见了那 98 块接下来要去哪儿。

书到这里就完了。它没能让你多赚一分钱，也没打算。它只在往后的每一次扫码、每一次看新闻、每一次签合同里，多给你留了一秒钟的明白。

有这一秒，就够了。

\begin{bridgebox}
\textbf{下一篇}全书讲了二十六章钱怎么跑，却始终没教你发财。这不是漏了，是故意的。后记二《我们为什么从头到尾不谈发财》，说这件事。
\end{bridgebox}

